\section{The attribute vocabulary}
\label{sec:attrs}

All parsing lives in \code{macros/teeny-macros/src/macros/tiled_kernel.rs}. Every
attribute below is host-only metadata: the macro strips it before re-emitting the
device and host signatures, and no attribute macro of that name is registered
anywhere.

\subsection{\texttt{\#[tile(...)]} --- one per real axis, outermost first}

A parameter carries one \code{#[tile(...)]} per axis it declares, in dim order
(dim~0 outermost, NCHW's \code{B} first). The axis's \emph{position} in that list is
its dim index, except that a \code{span} shifts every later index
(\S\ref{sec:span}).

{\small
\begin{xltabular}{\linewidth}{>{\raggedright\arraybackslash}p{3.6cm}>{\raggedright\arraybackslash}X}
\toprule
\textbf{Key} & \textbf{Meaning} \\
\midrule
\endhead
\code{extent = N} \emph{(required)} & The \code{N: i32} runtime parameter giving the axis's total extent. This is the axis's \emph{name} for propagation: two axes anywhere declaring the same extent are the same free variable. \\
\code{block = B} & A \code{const B: i32} generic, or a decimal literal, giving the tile size. Absent means one program per index. A literal (\code{block = 1}) records a kernel that steps the axis one element at a time and has no const to name. \\
\code{name = "B"} & The axis identity for \code{GridSpec} matching, and the suffix of the prelude's index binding (\code{tile_b}). It defaults to the extent's spelling; only \code{conv2d}'s \code{_B} needs it. \\
\code{dim = X|Y|Z} & The hardware grid dimension the axis decodes from. Defaults to \code{X}. \\
\code{window(stride, pad, kernel, output)} & The axis is read through a strided sliding window. \code{stride}/\code{pad}/\code{kernel} are const names or literals, \code{pad} is optional, and \code{output} names the \emph{output} axis the window resolves against. \\
\code{fill = zeros|neg_inf|pos_inf|one|<expr>} & The value of masked lanes: the reduction identity, or \code{constant_pad}'s runtime \code{value}. \\
\code{reduce} (flag) & This tensor is reduced over the axis. At most one per tensor. \\
\code{walk} (flag) & The axis is neither gridded nor loaded whole. A reduce pass or the map pass walks it in \code{block}-sized steps. \\
\code{span = N} & The axis covers \code{N} real dims, flattened row-major. Metadata only. Not allowed with \code{window}, \code{reduce}, or on a single-axis kernel (\S\ref{sec:span}). \\
\bottomrule
\end{xltabular}
}

Rules the macro enforces on route~1. Each is a compile error naming the offending
parameter.
\begin{itemize}
\item Either every \code{Tile} parameter declares \code{#[tile]} or none does. With
  none, the kernel falls back to the implicit \code{BLOCK_SIZE}/\code{n_elements}
  flat convention \emph{and no \code{tile_spec()} is generated}. Because that
  fallback is silent, \code{tests/test_tile_declarations.rs} parses the source to
  catch it.
\item There must be an \code{Out<Tile<..>>}, and it must carry a blocked axis. The
  output's axes are what the grid covers and what every input resolves against.
\item Every input axis must either match an output axis by \code{(extent, block,
  dim)}, be \code{reduce}, or carry a \code{window} whose \code{output} names an
  output axis. An input may declare a \emph{subset} of the output's axes; that is a
  broadcast operand (\S\ref{sec:prelude}).
\item A kernel cannot mix tagged \code{Tile} parameters with tagged raw pointers.
\end{itemize}

\subsection{Per-iteration reads: \texttt{\#[tile\_loop\_tile]} and \texttt{\#[tile\_loop\_scalar]}}

A loop-indexed operand cannot use the prelude's load-once read. Two
parameter attributes declare how it is read \emph{inside} the generated loop
instead. Both apply only to \code{In<Tile<..>>}.

\begin{description}[leftmargin=1.2em,style=nextline]
\item[{\code{#[tile_loop_tile(index = [coord = extent, ..], bounds = [i, ..])]}}]
  Reads a whole tile per iteration at loop-dependent coordinates. Each \code{index}
  entry is a full coordinate \emph{expression} and the layout extent, outermost
  first. The offset is the row-major fold of those pairs. \code{bounds} lists the
  positions whose coordinate gets a \code{ge(0) & lt(extent)} check (the padded
  axes). The coordinates are expressions rather than names on purpose: a windowed
  coordinate mixes a grid index, a loop index and two consts, and inferring which
  loop axis pairs with which window would be ambiguous. Requires exactly one
  blocked output axis.
\item[{\code{#[tile_loop_scalar(index = [idx = extent, ..])]}}]
  Reads one element per iteration at a loop-dependent flat index and broadcasts it
  to the carry's shape. This is conv2d's weight.
\end{description}

A loop-indexed operand still appears in the spec through its \code{#[tile]}
attributes. The two jobs are separate: \code{#[tile]} says what the operand
\emph{is}, and \code{#[tile_loop_tile]} says how to read it.

\subsection{Function attributes}

{\small
\begin{xltabular}{\linewidth}{>{\raggedright\arraybackslash}p{4.4cm}>{\raggedright\arraybackslash}X}
\toprule
\textbf{Attribute} & \textbf{Meaning} \\
\midrule
\endhead
\code{#[tile_loop(trip_count = [..])]} & An accumulation loop. \code{trip_count} is a list of \emph{names}, not a formula. It becomes \code{TileLoopSpec::trip_count_factors}. \\
\quad\code{axes = [name = extent, ..]} & The loop's own axes. Their extents multiply to the trip count, and the generated decode binds each name (innermost first). \\
\quad\code{count = <expr>} & An evaluable trip count, when there are no axes. \\
\quad\code{generate} (flag) & The macro emits the carry init, the loop around the author's body, and the store after it. \\
\code{#[tile_carry(acc = [shape])]} & A loop-carried accumulator. Each shape entry is a const name or a literal. \\
\quad\code{init = <expr>} & Replaces \code{zeros}. Max-pool uses $-\infty$. Requires \code{generate}. \\
\quad\code{finish = <expr>} & The epilogue: \code{let acc = <expr>;} between the loop and the store. Requires \code{generate}. \\
\code{#[tile_reduce_pass(over, read, into, acc, finish?, store?, fill?, merge?)]} & One walk over a \code{walk} axis accumulating into a scalar carry, then \code{finish} and an optional store. Passes run in declaration order. \code{merge} shares the previous pass's walk and read. \\
\code{#[tile_grid(order = [..])]} & The real grid decode order when it differs from the output's dim order, or a subset when a loop covers an axis. \\
\code{#[tile_grid(swizzled)]} & No \code{GridSpec} at all: the decode is not mixed-radix (matmul's \code{GROUP_M}). \\
\bottomrule
\end{xltabular}
}

Errors the loop attributes guard against:
\begin{itemize}
\item a carry without a loop;
\item a loop without a carry (``an independent walk belongs in the grid'');
\item \code{generate} with neither \code{axes} nor \code{count};
\item \code{finish}/\code{init} without \code{generate} (otherwise the expression is silently dropped);
\item \code{generate} with more than one carry or more than one \code{Out<Tile>} (flash attention's three carries need a carry-to-output pairing that does not exist yet);
\item a \code{merge} pass whose \code{read}/\code{over} differs from its leader's.
\end{itemize}

\subsection{\texttt{span}: one declared axis, several real dims}
\label{sec:span}

New in PR \#103. A kernel may address several real dims as one row-major
flattened axis. \code{batch_norm_2d_nchw_forward_inference} walks an NCHW tensor
as \code{[B, C, HW]}. \code{span = 2} tells the spec the truth about the tensor
without touching the body:

\begin{lstlisting}
#[tile(extent = B, dim = Y)]
#[tile(extent = C)]
#[tile(block = BLOCK_HW, extent = HW, span = 2)]   // dims 2 and 3
x: In<Tile<T, D>>,
\end{lstlisting}

Only the metadata moves. The spec's rank is $\sum \text{span}$ (4 here), the
binding's \code{dims} are \code{[2, 3]}, and \code{reduction_axis} is now a real dim
index rather than a declaration position. The prelude still decodes and addresses
one \code{HW} axis. The combinations that would need a meaning the spec types cannot
carry are refused:
\begin{itemize}
\item \code{window}: a window relates one input dim to one output dim.
\item \code{reduce}: \code{reduction_axis} holds a single index.
\item A single-axis kernel: its \code{tile_spec(rank)} already flattens every dim.
\end{itemize}
